\documentclass[a4paper,12pt,oneside]{article}

\usepackage[T2A]{fontenc}
\usepackage[utf8]{inputenc}
\usepackage[russian]{babel}

\usepackage{microtype}

\usepackage{amsmath,amssymb,mathtools}

\usepackage{geometry}
\geometry{top=25mm,left=25mm,right=20mm,bottom=25mm}

\usepackage{enumitem}
\setlist[enumerate]{leftmargin=*, itemsep=0.55\baselineskip, topsep=0.35\baselineskip}
\setlist[itemize]{leftmargin=*, itemsep=0.25\baselineskip, topsep=0.25\baselineskip}

\setlength{\parindent}{1.25em}
\setlength{\parskip}{0pt}

\newcommand{\Task}[1]{\section*{\normalfont\bfseries #1}}

\begin{document}

\begin{center}
    {\Large\bfseries Домашняя работа №4. Кодирование}\\[2mm]
\end{center}
\vspace{2mm}
\hrule
\vspace{4mm}

\Task{Закодируйте своё ФИО без учёта пробелов и регистра}

\begin{enumerate}[label=\arabic*.]
\item Равномерным кодом.
\item Кодом Хаффмана. Построить дерево.
\item Кодом Шеннона--Фано. Заполнить таблицу.
\item Арифметическим кодом. Заполнить таблицы с левыми и правыми границами. Допускается использовать программу (код прилагается к работе), однако все отрезки должны быть приведены в работе. Построить бинарный код.
\item Оцените энтропию.
\item Сравните средние длины кодовых слов и энтропию.
\item Результат арифметического кодирования закодируйте кодом Хэмминга. Оцените \textit{Rate}.
\end{enumerate}

\bigskip

\textbf{Важно:} для каждого разделимого кода продемонстрировать таблицу кодовых слов.

\end{document}
